\documentclass{article}

\usepackage{arxiv}

\usepackage[utf8]{inputenc}
\usepackage[T1]{fontenc}
\usepackage{hyperref}
\usepackage{url}
\usepackage{booktabs}
\usepackage{amsmath, amssymb, amsthm}
\usepackage{amsfonts}
\usepackage{bm}
\usepackage{nicefrac}
\usepackage{microtype}
\usepackage{graphicx}
\usepackage{xcolor}

\pagecolor[rgb]{0.05,0.05,0.05} %black - background color
\color[rgb]{0.9,0.9,0.9} %grey - text color


\newcommand{\synfig}[1]{notebooks/figures/synthetic/gene_expression_synthetic/imaging_synthetic/#1}
\newcommand{\kircfig}[1]{notebooks/figures/tcga_kirc/gene_expression/tumor_radimagenet/#1}

\newtheorem{proposition}{Proposition}
\newtheorem{definition}{Definition}
\newtheorem{corollary}{Corollary}

\newcommand{\R}{\mathbb{R}}
\newcommand{\E}{\mathbb{E}}
\newcommand{\N}{\mathcal{N}}
\newcommand{\Cov}{\operatorname{Cov}}
\newcommand{\Var}{\operatorname{Var}}
\newcommand{\rank}{\operatorname{rank}}
\newcommand{\tr}{\operatorname{tr}}
\newcommand{\Sg}{\Sigma_{G}}
\newcommand{\Sx}{\Sigma_{X}}
\newcommand{\Sgx}{\Sigma_{GX}}
\newcommand{\Sxg}{\Sigma_{XG}}
\newcommand{\Sgcx}{\Sigma_{G\mid X}}
\newcommand{\T}{^{\top}}

\title{Image-Identifiable Genomic Subspaces:\\
A Linear--Gaussian Model of Radiogenomic Recoverability}

\author{
  Joseph Rich \\
  Division of Biology and Biological Engineering\\
  California Institute of Technology \\
  Keck School of Medicine, University of Southern California \\
  \And
  Lior Pachter\thanks{Address correspondence to \texttt{lpachter@caltech.edu}.} \\
  Division of Biology and Biological Engineering\\
  Department of Computing and Mathematical Sciences\\
  California Institute of Technology \\
}

\date{}

\renewcommand{\shorttitle}{Image-Identifiable Genomic Subspaces}

\hypersetup{
pdftitle={Image-Identifiable Genomic Subspaces: A Linear-Gaussian Model of Radiogenomic Recoverability},
pdfauthor={Joseph Rich, Lior Pachter},
pdfkeywords={radiogenomics, canonical correlation analysis, mutual information, recoverability, linear-Gaussian model},
}

\begin{document}
\maketitle

\begin{abstract}
We develop a linear--Gaussian generative model linking patient-level genomic
profiles $g \in \R^{p}$ to imaging phenotypes $x \in \R^{d}$, and use it to ask
a sharper question than ``does imaging predict gene $j$?'' Instead we ask
\emph{which low-dimensional genomic subspaces are image-identifiable}. Under the
model every quantity of interest is available in closed form: the Bayes-optimal
estimator $\E[g \mid x]$, the posterior covariance $\Sgcx$, the mutual
information $I(g;x)$, and a per-direction \emph{recoverability score}
$R(v) \in [0,1]$. We show that the directions of maximal recoverability solve a
generalized eigenproblem that coincides exactly with canonical correlation
analysis (CCA) between the two modalities, and that the recoverability of the
$i$-th canonical genomic direction equals the squared canonical correlation
$\rho_i^2$. We then address estimation in the regime $p \gg n$ via regularized
covariances, reduced-rank regression, and the probabilistic-CCA latent-variable
form, and describe cross-validated and permutation-based procedures for honest
inference. The companion notebook
\texttt{notebooks/radiogenomic\_recoverability.ipynb} implements the full
pipeline on a pair of \texttt{AnnData} objects, with a synthetic-data fallback
that has known ground truth.
\end{abstract}

\keywords{radiogenomics \and canonical correlation analysis \and mutual information \and recoverability \and linear--Gaussian model}

\section{Setup and notation}
\label{sec:setup}

We observe $n$ patients. For patient $i$ let $g_i \in \R^{p}$ be a vector of
genomic features (e.g.\ expression, somatic mutation burden, copy-number scores)
and $x_i \in \R^{d}$ a vector of imaging phenotypes (radiomic descriptors, deep
features, or structured reads). Stack the observations into data matrices
$\mathbf{G} \in \R^{n \times p}$ and $\mathbf{X} \in \R^{n \times d}$, with row
$i$ equal to $g_i\T$ and $x_i\T$ respectively. Random variables are written in
uppercase ($G, X$), realizations in lowercase ($g, x$). Throughout we assume
both modalities have been centered, so all means are zero; recentring is handled
by the estimators of Section~\ref{sec:estimation}.

We use $\Sg \in \R^{p \times p}$, $\Sx \in \R^{d \times d}$ for the marginal
covariances, $\Sgx = \Cov(G, X) \in \R^{p \times d}$ for the cross-covariance,
and $\Sxg = \Sgx\T$. For a symmetric positive semidefinite $S$ we write $S^{1/2}$ for
its symmetric square root and $S^{-1/2}$ for the inverse thereof.

\section{The linear--Gaussian generative model}
\label{sec:model}

The genomic state is drawn from a Gaussian prior and the imaging phenotype is a
noisy linear readout of it:
\begin{align}
G &\sim \N(0, \Sg), \label{eq:prior}\\
X \mid G &\sim \N(A G,\ \sigma^2 I_d), \qquad A \in \R^{d \times p}. \label{eq:likelihood}
\end{align}
The matrix $A$ encodes how genomic variation is expressed in the image; the
isotropic term $\sigma^2 I_d$ collects everything in the image that genomics does
not drive --- biological variation outside the measured panel, acquisition and
reconstruction noise, and feature-extraction error. Equivalently,
\begin{equation}
X = A G + \varepsilon, \qquad \varepsilon \sim \N(0, \sigma^2 I_d), \quad \varepsilon \perp G.
\label{eq:structural}
\end{equation}
Equations~\eqref{eq:prior}--\eqref{eq:structural} are deliberately the simplest
model with a non-trivial answer. Their value is that \emph{every} downstream
quantity is closed-form, so the model functions as an exact upper bound: no
estimator --- linear or nonlinear --- can recover genomics from imaging better
than the posterior derived below (under these assumptions). Section~\ref{sec:estimation} relaxes the
isotropic-noise and known-parameter assumptions for practical fitting.

\section{Marginal and joint laws}
\label{sec:joint}

\begin{proposition}[Joint Gaussian law]
\label{prop:joint}
Under \eqref{eq:prior}--\eqref{eq:structural},
\begin{equation}
\begin{pmatrix}
G\\
X
\end{pmatrix}
\sim
\N\!\left(
\bf{0},
\begin{pmatrix}
\Sg & \Sg A^\top\\
A\Sg & A\Sg A^\top + \sigma^2 I_d
\end{pmatrix}
\right).
\end{equation}
In particular,
\begin{equation}
X \sim \N(0,\Sx),
\qquad
\Sx = A\Sg A^\top + \sigma^2 I_d,
\label{eq:sigmax}
\end{equation}
and the cross-covariance is
\begin{equation}
\Sgx = \Cov(G,X)=\Sg A^\top.
\end{equation}
\end{proposition}

\begin{proof}
Since
\[
X = AG + \varepsilon,
\]
where
\[
G \sim \N(0,\Sg),
\qquad
\varepsilon \sim \N(0,\sigma^2 I_d),
\]
and $G$ and $\varepsilon$ are independent, the stacked vector
\[
\begin{pmatrix}
G\\
X
\end{pmatrix}
=
\begin{pmatrix}
G\\
AG+\varepsilon
\end{pmatrix}
\]
is jointly Gaussian, as it is an affine transformation of independent Gaussian variables.

It therefore suffices to compute its mean and covariance. Since both $G$ and $\varepsilon$ are centered,
\[
\E[G]=0,
\qquad
\E[\varepsilon]=0,
\]
we obtain
\[
\E[X]
=
A\E[G]+\E[\varepsilon]
=
0,
\]
and hence
\[
\E
\begin{pmatrix}
G\\
X
\end{pmatrix}
=
0.
\]

Next, the covariance matrix has block form
\[
\Cov
\begin{pmatrix}
G\\
X
\end{pmatrix}
=
\begin{pmatrix}
\Cov(G,G) & \Cov(G,X)\\
\Cov(X,G) & \Cov(X,X)
\end{pmatrix}.
\]

The upper-left block is simply
\[
\Cov(G,G)=\Sg.
\]

For the upper-right block, using $X=AG+\varepsilon$,
\[
\Cov(G,X)
=
\Cov(G,AG+\varepsilon).
\]
By linearity of covariance,
\[
\Cov(G,X)
=
\Cov(G,AG)+\Cov(G,\varepsilon).
\]
Since $G$ and $\varepsilon$ are independent,
\[
\Cov(G,\varepsilon)=0,
\]
and therefore
\[
\Cov(G,X)
=
\Cov(G,AG).
\]
Using the identity
\[
\Cov(U,BV)=\Cov(U,V)B^\top,
\]
we obtain
\[
\Cov(G,AG)
=
\Sg A^\top.
\]
Hence
\[
\Sgx=\Cov(G,X)=\Sg A^\top.
\]

The lower-left block is the transpose:
\[
\Cov(X,G)
=
\Cov(G,X)^\top
=
A\Sg.
\]

Finally, for the lower-right block,
\[
\Cov(X,X)
=
\Cov(AG+\varepsilon).
\]
Expanding the covariance,
\[
\Cov(X,X)
=
\Cov(AG)+\Cov(\varepsilon),
\]
since the cross terms vanish by independence. Using
\[
\Cov(AG)=A\Sg A^\top,
\]
together with
\[
\Cov(\varepsilon)=\sigma^2 I_d,
\]
gives
\[
\Cov(X,X)
=
A\Sg A^\top+\sigma^2 I_d.
\]
Therefore
\[
\Sx=A\Sg A^\top+\sigma^2 I_d.
\]

Collecting the covariance blocks yields
\[
\Cov
\begin{pmatrix}
G\\
X
\end{pmatrix}
=
\begin{pmatrix}
\Sg & \Sg A^\top\\
A\Sg & A\Sg A^\top+\sigma^2 I_d
\end{pmatrix},
\]
which proves the claim.
\end{proof}

% \begin{proof}
% $G$ and $\varepsilon$ are jointly Gaussian and independent, and
% $(G, X)$ is a linear image of $(G, \varepsilon)$, hence jointly Gaussian. The
% blocks follow from $\Cov(G,X) = \Cov(G, AG + \varepsilon) = \Sg A\T$ and
% $\Var(X) = A \Sg A\T + \sigma^2 I_d$.
% \end{proof}

Equation~\eqref{eq:sigmax} is the model's first interpretive statement: imaging
covariance decomposes additively into a \emph{genomics-driven} component
$A \Sg A\T$ and an \emph{irreducible} component $\sigma^2 I_d$. The cross-covariance
$\Sgx = \Sg A\T$ is the central estimable object --- it names which genomic
directions co-vary with imaging at all.

\section{Bayes-optimal recovery of genomics from imaging}
\label{sec:posterior}

\begin{proposition}[Posterior]
\label{prop:posterior}
The posterior of genomics given imaging is Gaussian, $G \mid X \sim \N(\hat g(X), \Sgcx)$, with
\begin{align}
\hat g(X) &= \E[G \mid X] = \Sg A\T \big(A \Sg A\T + \sigma^2 I_d\big)^{-1} X \;=\; B X, \label{eq:postmean}\\
\Sgcx &= \Sg - \Sg A\T \big(A \Sg A\T + \sigma^2 I_d\big)^{-1} A \Sg. \label{eq:postcov}
\end{align}
Equivalently, in precision form,
\begin{equation}
\Sgcx^{-1} = \Sg^{-1} + \sigma^{-2} A\T A,
\qquad
B = \sigma^{-2}\, \Sgcx\, A\T .
\label{eq:precision}
\end{equation}
\end{proposition}

\begin{proof}
Equations~\eqref{eq:postmean}--\eqref{eq:postcov} are the standard Gaussian
conditioning formulas applied to Proposition~\ref{prop:joint}. The precision
form follows from the Woodbury identity:
$\big(\Sg^{-1} + \sigma^{-2} A\T A\big)^{-1} = \Sg - \Sg A\T(A\Sg A\T + \sigma^2 I_d)^{-1} A \Sg$,
and then $B = \Sg A\T \Sx^{-1} = \sigma^{-2}\Sgcx A\T$ by the push-through
identity $\Sg A\T(A\Sg A\T+\sigma^2 I)^{-1} = (\Sg^{-1}+\sigma^{-2}A\T A)^{-1}\sigma^{-2}A\T$.
\end{proof}

The estimator $\hat g(X) = BX$ is the Bayes-optimal predictor under squared-error
loss; because the model is jointly Gaussian it is also the optimal predictor over
\emph{all} measurable functions of $X$, not merely linear ones. Thus $\Sgcx$ is a
hard floor on residual genomic uncertainty: \textbf{no algorithm can do better}.
The precision form~\eqref{eq:precision} makes the accounting transparent ---
imaging \emph{adds} information $\sigma^{-2} A\T A$ to the prior precision
$\Sg^{-1}$.

\section{Mutual information and the SNR decomposition}
\label{sec:mi}

\begin{proposition}[Mutual information]
\label{prop:mi}
The mutual information between the two modalities is
\begin{equation}
I(G; X)
= \tfrac12 \log \frac{\det \Sg}{\det \Sgcx}
= \tfrac12 \log \det\!\big(I_d + \sigma^{-2} A \Sg A\T\big).
\label{eq:mi}
\end{equation}
Let $\lambda_1 \ge \cdots \ge \lambda_r > 0$ be the nonzero eigenvalues of
$A \Sg A\T$. Then
\begin{equation}
I(G; X) = \tfrac12 \sum_{i=1}^{r} \log\!\left(1 + \frac{\lambda_i}{\sigma^2}\right).
\label{eq:mi-diag}
\end{equation}
\end{proposition}

\begin{proof}
For jointly Gaussian $(G,X)$,
$I(G;X) = \tfrac12 \log\big(\det\Sg / \det\Sgcx\big)$. Substituting
\eqref{eq:precision}, $\det \Sgcx = \det\Sg \,/ \det(I_p + \sigma^{-2}\Sg^{1/2}A\T A \Sg^{1/2})$,
and Sylvester's determinant identity moves $A$ to the other side to give
$\det(I_d + \sigma^{-2} A\Sg A\T)$. Diagonalizing $A\Sg A\T$ yields
\eqref{eq:mi-diag}.
\end{proof}

Equation~\eqref{eq:mi-diag} is the cleanest reading of the model: each
image-visible genomic direction contributes
$\tfrac12 \log(1 + \mathrm{SNR}_i)$ bits (nats), where
$\mathrm{SNR}_i = \lambda_i / \sigma^2$ is the signal-to-noise ratio of that
direction. Directions whose genomic signal is small relative to imaging noise
contribute essentially nothing, regardless of how many genes load on them.
Figure~\ref{fig:synthetic-mi} illustrates the decomposition
\eqref{eq:mi-diag} on the synthetic instance introduced below: the per-direction
bit contribution falls off rapidly after the first handful of canonical
components, and the cumulative total saturates at $I(G;X)\approx 6.45$ bits
out of the analytic upper bound implied by the model.

\begin{figure}[h]
\centering
\includegraphics[width=0.75\linewidth]{\synfig{mi_per_direction.pdf}}
\caption{Per-direction contribution to $I(G;X)$ on a synthetic linear--Gaussian
instance ($n=800$, $p=80$, $d=40$). Each bar is
$\tfrac12 \log(1 + d_i^2)$ in bits; the dashed line is the cumulative total.}
\label{fig:synthetic-mi}
\end{figure}

\section{Canonical coordinates and the recoverability spectrum}
\label{sec:canonical}

The eigenvalues $\lambda_i$ are most interpretable after whitening. Define the
whitened genomic variable $U = \Sg^{-1/2} G \sim \N(0, I_p)$ and the
\emph{whitened transfer operator}
\begin{equation}
M := \sigma^{-1} A\, \Sg^{1/2} \in \R^{d \times p},
\qquad
M = \sum_{i} d_i\, u_i v_i\T \quad\text{(SVD)},
\label{eq:M}
\end{equation}
with left/right singular vectors $u_i \in \R^d$, $v_i \in \R^p$ and singular
values $d_1 \ge d_2 \ge \cdots \ge 0$.

\begin{proposition}[Canonical decomposition of recoverability]
\label{prop:canonical}
With $M$ as in \eqref{eq:M}, define the genomic canonical direction
$a_i := \Sg^{-1/2} v_i$ and the canonical score $S_i := a_i\T G = v_i\T U$. Then:
\begin{enumerate}
\item[(i)] $S_i$ has prior variance $1$ and posterior variance
$\Var(S_i \mid X) = (1 + d_i^2)^{-1}$;
\item[(ii)] the canonical correlation between $S_i$ and imaging is
$\rho_i = d_i / \sqrt{1 + d_i^2}$, equivalently $d_i^2 = \rho_i^2/(1-\rho_i^2)$;
\item[(iii)] the eigenvalues of Section~\ref{sec:mi} satisfy
$\lambda_i / \sigma^2 = d_i^2$, so
\begin{equation}
I(G;X) = \tfrac12 \sum_i \log(1 + d_i^2) = -\tfrac12 \sum_i \log(1 - \rho_i^2).
\end{equation}
\end{enumerate}
\end{proposition}

\begin{proof}
In whitened coordinates the likelihood is $X \mid U \sim \N(\sigma M U, \sigma^2 I_d)$,
since $A G = A \Sg^{1/2} U = \sigma M U$. The posterior precision of $U$ is
$I_p + \sigma^{-2}(\sigma M)\T(\sigma M) = I_p + M\T M = V (I_p + D^2) V\T$,
so $\Var(U \mid X) = V(I_p + D^2)^{-1} V\T$ and
$\Var(S_i \mid X) = v_i\T \Var(U\mid X) v_i = (1+d_i^2)^{-1}$, giving (i).
For (ii), $\Cov(U, X) = \E[U(\sigma M U + \varepsilon)\T] = \sigma M\T$, so
$\Cov(S_i, X) = \sigma v_i\T M\T = \sigma d_i u_i\T$. The matched imaging variate
$T_i := u_i\T X$ has $\Var(T_i) = u_i\T \Sx u_i = \sigma^2(1 + d_i^2)$ (using
$\Sx = \sigma^2(MM\T + I_d)$) and $\Cov(S_i, T_i) = \sigma d_i$, hence
$\rho_i = \mathrm{corr}(S_i, T_i) = d_i/\sqrt{1+d_i^2}$. Part (iii) is immediate
from $A\Sg A\T = \sigma^2 M M\T$ and $\det(I_d + MM\T) = \prod_i(1+d_i^2)$.
\end{proof}

Part (i) gives the headline identity. Define the recoverability of a genomic
score as one minus the fraction of its variance that survives observing imaging
(formalized in Section~\ref{sec:scores}); then
\begin{equation}
\boxed{\;R(a_i) = 1 - \Var(S_i \mid X) = 1 - \frac{1}{1+d_i^2} = \frac{d_i^2}{1+d_i^2} = \rho_i^2.\;}
\label{eq:R-rho}
\end{equation}
\textbf{The recoverability of the $i$-th canonical genomic direction is exactly
the squared canonical correlation.} This is what reduces the whole analysis to
CCA between $\mathbf{G}$ and $\mathbf{X}$ (Section~\ref{sec:estimation}), and it
means $A$ and $\sigma^2$ never have to be estimated individually --- only the
three second moments $\Sg, \Sx, \Sgx$ enter.

\section{Recoverability of a named genomic direction}
\label{sec:scores}

In practice one often cares about a specific, biologically named genomic score
$Y = v\T G$ for a fixed $v \in \R^p$ --- a pathway signature, a polygenic score,
a single gene's expression ($v = e_j$).

\begin{definition}[Recoverability score]
For $v \in \R^p$ with $v\T \Sg v > 0$,
\begin{equation}
R(v) := 1 - \frac{v\T \Sgcx\, v}{v\T \Sg\, v}
= \frac{v\T (\Sg - \Sgcx)\, v}{v\T \Sg\, v}
\;\in\; [0, 1].
\label{eq:Rv}
\end{equation}
\end{definition}

$R(v) = 0$ means imaging carries no information about $Y = v\T G$;
$R(v) = 1$ means imaging determines it exactly. The score is scale-invariant in
$v$ and, by the Gaussian conditional-variance identity, equals the population
$R^2$ of the Bayes-optimal predictor of $Y$ from $X$:
$R(v) = 1 - \E[\Var(Y \mid X)] / \Var(Y)$. The matrix at the core of
\eqref{eq:Rv} has the convenient moment-only form
\begin{equation}
\Sg - \Sgcx = \Sg A\T \Sx^{-1} A \Sg = \Sgx\, \Sx^{-1}\, \Sxg .
\label{eq:explained}
\end{equation}

\section{The recoverability eigenproblem}
\label{sec:eigenproblem}

The natural project-level question --- \emph{which} genomic directions does
imaging see best --- is the variational problem
\begin{equation}
\max_{v \neq 0}\;
R(v)
= \max_{v \neq 0}\;
\frac{v\T \big(\Sgx \Sx^{-1} \Sxg\big) v}{v\T \Sg\, v}.
\label{eq:varproblem}
\end{equation}

\begin{proposition}[Recoverability eigenproblem $=$ CCA]
\label{prop:eig}
The stationary points of \eqref{eq:varproblem} are the solutions of the
generalized eigenproblem
\begin{equation}
\Sgx\, \Sx^{-1}\, \Sxg\; v = \mu\, \Sg\, v .
\label{eq:gevp}
\end{equation}
Its eigenvalues are $\mu_i = \rho_i^2 = d_i^2/(1+d_i^2)$ and its eigenvectors are
the canonical genomic directions $a_i$ of Proposition~\ref{prop:canonical},
$\Sg$-orthonormal ($a_i\T \Sg a_j = \delta_{ij}$). Thus the ordered solutions of
\eqref{eq:gevp} are precisely the canonical correlation directions between $G$
and $X$, and $\max_v R(v) = \rho_1^2$.
\end{proposition}

\begin{proof}
\eqref{eq:gevp} is the first-order condition of the Rayleigh quotient
\eqref{eq:varproblem}. Substitute $v = \Sg^{-1/2} w$ and $\Sgx = \Sg A\T$; the
generalized eigenproblem becomes the ordinary one
$\Sg^{1/2} A\T \Sx^{-1} A \Sg^{1/2}\, w = \mu\, w$. Now $A \Sg^{1/2} = \sigma M$
and $\Sx = \sigma^2(MM\T + I_d)$, so
\begin{equation*}
\Sg^{1/2} A\T \Sx^{-1} A \Sg^{1/2}
= \sigma M\T \cdot \sigma^{-2}(MM\T + I_d)^{-1} \cdot \sigma M
= M\T (MM\T + I_d)^{-1} M .
\end{equation*}
With the SVD $M = U D V\T$ this equals $V\, D^2 (D^2 + I)^{-1}\, V\T$, whose
eigenvalues are $d_i^2/(1+d_i^2) = \rho_i^2$ with eigenvectors $v_i$. Hence
$w = v_i$ and $v = \Sg^{-1/2} v_i = a_i$; the $\Sg$-orthonormality
$a_i\T \Sg a_j = v_i\T v_j = \delta_{ij}$ is immediate.
\end{proof}

So the deliverable of the analysis is a \emph{recoverability spectrum}
$\rho_1^2 \ge \rho_2^2 \ge \cdots$ together with the genomic directions $a_i$
that achieve it. Reading the loadings of $a_1, a_2, \dots$ back onto genes or
pathways answers ``which biology is image-identifiable'': the leading directions
typically align with coarse, spatially expressed programs --- proliferation,
hypoxia, angiogenesis, necrosis, immune infiltration, tumor burden, stromal
remodeling --- rather than with individual transcripts.

\paragraph{Illustration on synthetic data.}
Figure~\ref{fig:synthetic-spectrum} shows the recoverability spectrum on a
synthetic radiogenomic instance generated from the linear--Gaussian model of
Section~\ref{sec:model} ($n=800$, $p=80$, $d=40$, latent rank $k=5$ with
prescribed $\{\rho_i^\star\}$). The estimated $\hat\rho_i^2$ closely match
the planted values for the identifiable directions and decay toward the
high-dimensional bulk thereafter, validating
Proposition~\ref{prop:eig}. Figure~\ref{fig:synthetic-truth} reads the same
estimator against the analytic ground truth derived from
Propositions~\ref{prop:posterior}--\ref{prop:canonical}: the recovered
posterior mean $\hat B X$, posterior-variance retention, and direction
alignment all track the closed-form predictions.

\begin{figure}[h]
\centering
\includegraphics[width=0.85\linewidth]{\synfig{recoverability_spectrum.pdf}}
\caption{Recoverability spectrum on a synthetic linear--Gaussian instance
($n=800$, $p=80$, $d=40$). The leading $\hat\rho_i^2$ recovers the planted
canonical structure; later components fall to the high-dimensional bulk.}
\label{fig:synthetic-spectrum}
\end{figure}

\begin{figure}[h]
\centering
\includegraphics[width=0.85\linewidth]{\synfig{synthetic_truth_validation.pdf}}
\caption{Empirical-vs-analytic check on the same synthetic instance.
Estimated canonical correlations, recoverability per direction, and
imaging variance attributable to genomics all match the closed-form
quantities from Propositions~\ref{prop:posterior}--\ref{prop:canonical}.}
\label{fig:synthetic-truth}
\end{figure}

\section{Rank limits and identifiability}
\label{sec:rank}

\begin{proposition}[Rank limit]
\label{prop:rank}
The recoverable signal matrix has rank
\begin{equation}
\rank(\Sg - \Sgcx) = \rank(A \Sg A\T)
\le \min\!\big(d,\ p,\ \rank A,\ \rank \Sg\big).
\end{equation}
Consequently at most $\min(d, p, \rank A, \rank\Sg)$ canonical correlations are
nonzero, and the image-identifiable genomic subspace
$\operatorname{span}\{a_i : \rho_i > 0\}$ has dimension bounded by the same
quantity.
\end{proposition}

This is the conceptual core. Even with $p = 20{,}000$ measured genes, imaging can
expose only a projection of genomics whose dimension is capped by the rank of the
transfer map $A$ --- if a CT phenotype reflects, say, $20$ macroscopic
biological axes, then no model recovers $20{,}000$ independent molecular
variables from it. The recoverable dimension is a property of the
\emph{biology--imaging channel}, not of the estimator. A data-processing remark:
imaging features are themselves a deterministic, lossy function
$x = \phi(\text{image})$; any such post-processing can only shrink
$\rank(\Sgx)$ and lower each $\rho_i$, never raise them.

\section{Estimation from finite samples}
\label{sec:estimation}

In applications $p$ and $d$ may rival or exceed $n$, the noise is not isotropic,
and $A, \Sg, \sigma^2$ are unknown. Because Propositions~\ref{prop:canonical}
and~\ref{prop:eig} require only $\Sg, \Sx, \Sgx$, we estimate those directly.

\subsection{Regularized covariances}
\label{sec:estimation-cov}
Replace sample covariances by shrinkage estimators (e.g.\ Ledoit--Wolf),
$\hat\Sg = (1-\alpha_G)\,S_G + \alpha_G\, \nu_G I_p$ and likewise $\hat\Sx$, with
$S_G = n^{-1}\mathbf{G}\T\mathbf{G}$ the centered sample covariance; the
cross-covariance is $\hat\Sgx = n^{-1}\mathbf{G}\T\mathbf{X}$. Shrinkage keeps
$\hat\Sg, \hat\Sx$ invertible and well-conditioned, which is required for the
whitening in \eqref{eq:M}. When $p$ is very large it is practical to first
project genomics onto its top principal components (or onto curated pathway
scores), reducing $p$ to a manageable $p' \ll n$ before estimation; the rank
limit of Proposition~\ref{prop:rank} means little is lost provided $p'$ exceeds
the recoverable dimension.

\subsection{Reduced-rank and ridge regression}
\label{sec:estimation-rrr}
The transfer map, if needed explicitly, is the multivariate regression of $X$ on
$G$. Ridge gives
$\hat A\T = (\mathbf{G}\T\mathbf{G} + \gamma I_p)^{-1}\mathbf{G}\T\mathbf{X}$,
and $\hat\sigma^2$ follows from the residual covariance. Imposing
$\rank(A) \le k$ (reduced-rank regression) bakes Proposition~\ref{prop:rank}
into the fit and is equivalent, up to regularization, to truncating the SVD of
$M$ at $k$ components.

\subsection{Regularized CCA}
\label{sec:estimation-cca}
The estimator used in the companion notebook solves \eqref{eq:gevp} with the
regularized covariances of Section~\ref{sec:estimation-cov}: form
$\hat M = \hat\Sg^{-1/2} \hat\Sgx \hat\Sx^{-1/2}$, take its SVD
$\hat M = \hat U \hat D \hat V\T$, and read off
$\hat\rho_i = \hat d_i$ (clipped to $[0,1)$),
$\hat R_i = \hat\rho_i^2$, genomic directions $\hat a_i = \hat\Sg^{-1/2}\hat u_i$,
imaging directions $\hat b_i = \hat\Sx^{-1/2}\hat v_i$. The Bayes-optimal
predictor and posterior covariance follow from
\eqref{eq:postmean}--\eqref{eq:postcov} with hats throughout. This is exactly
regularized CCA --- the recoverability program and CCA are the same computation.

\subsection{Probabilistic CCA: the latent-variable view}
\label{sec:estimation-pcca}
The rank limit motivates an explicit low-dimensional generative form. With a
shared latent $Z \in \R^{k}$,
\begin{equation}
Z \sim \N(0, I_k), \qquad
G = W_G Z + \mu_G + \epsilon_G, \qquad
X = W_X Z + \mu_X + \epsilon_X,
\label{eq:pcca}
\end{equation}
$\epsilon_G \sim \N(0, \Psi_G)$, $\epsilon_X \sim \N(0, \Psi_X)$ independent
(probabilistic CCA; Bach \& Jordan, 2005). Here $k$ \emph{is} the dimension of
the image-identifiable genomic subspace, $Z$ is the shared biological state, and
$\epsilon_G, \epsilon_X$ are modality-private variation. The implied moments are
$\Sg = W_G W_G\T + \Psi_G$, $\Sx = W_X W_X\T + \Psi_X$,
$\Sgx = W_G W_X\T$; the maximum-likelihood loadings are the canonical directions
scaled by the canonical correlations, recovering Section~\ref{sec:estimation-cca}.
The model is fit by EM, which also yields posterior latent estimates
$\E[Z \mid x]$ and naturally accommodates anisotropic (per-feature) noise, unlike
the isotropic $\sigma^2 I_d$ of the base model.

\subsection{Sample-size regimes and dimensional collapse}
\label{sec:estimation-regimes}
Equations~\eqref{eq:postmean}--\eqref{eq:gevp} are population identities; the
estimators of Sections~\ref{sec:estimation-cov}--\ref{sec:estimation-pcca}
inherit a finite-sample bias whose magnitude is governed entirely by the
aspect ratios $p/n$ and $d/n$. The bias is not benign, and below a critical
sample size the recoverability spectrum breaks down in a specific, predictable
way.

\paragraph{Three regimes.} Let $p^\star, d^\star$ be the working dimensions
after any pre-reduction (e.g.\ PCA), and write $\kappa_G = p^\star/n$,
$\kappa_X = d^\star/n$.
\begin{enumerate}
\item \emph{Classical} ($\kappa_G + \kappa_X \ll 1$). Sample CCA is consistent;
$\hat\rho_i \to \rho_i$ at the usual $n^{-1/2}$ rate and the in-sample spectrum
is a faithful estimate of the population one.
\item \emph{High-dimensional} ($\kappa_G + \kappa_X \lesssim 1$, both bounded
away from $1$). Under the null $\Sgx = 0$, the squared sample canonical
correlations follow a Wachter-type law on
$[(\sqrt{\kappa_G(1-\kappa_X)} - \sqrt{\kappa_X(1-\kappa_G)})^2,
 (\sqrt{\kappa_G(1-\kappa_X)} + \sqrt{\kappa_X(1-\kappa_G)})^2]$
(Wachter, 1980; Yang \& Pan, 2015); the leading sample value already sits well
above zero by geometry. Estimates of $\rho_i$ are systematically inflated and
ordinary CCA must be replaced by its regularized form, with shrinkage
parameters tuned by cross-validation.
\item \emph{Saturating} ($p^\star + d^\star \ge n$). The rows of $\mathbf{G}$
and $\mathbf{X}$ live in subspaces whose intersection has dimension at least
$p^\star + d^\star - n + 1$; consequently
$\hat\rho_1 = \cdots = \hat\rho_q = 1$ for $q = \min(p^\star, d^\star) - (n -
\max(p^\star, d^\star))_+$ \emph{regardless of any underlying association}.
The empirical mutual information diverges, the in-sample
$\tr(\hat\Sgx \hat\Sx^{-1} \hat\Sxg)/\tr(\hat\Sx)$ approaches the rank fraction
of the genomic span, and the permutation null collapses onto the observed
value (its $p$-values concentrate near $1$). Nothing in the spectrum is
estimable in this regime without shrinkage \emph{plus} dimension reduction;
shrinkage alone keeps the matrices invertible but does not restore
identifiability.
\end{enumerate}

\paragraph{Quantitative warning signs.} Three diagnostics, jointly, separate
real signal from dimensional artefact:
\begin{itemize}
\item the gap between in-sample $\hat R_i$ and the $K$-fold value $\hat
R_i^{\mathrm{cv}}$; under regime~3 this gap approaches $\hat R_i - 0$;
\item the permutation null for $\hat\rho_1$: if the observed value lies inside
the bulk, the result is consistent with chance;
\item the cross-validated $R^2$ of the leading canonical score $a_1\T g$
predicted from $x$ on a held-out split; in regime~3 it is typically large and
\emph{negative} (worse than predicting the mean), a behaviour seen empirically
on TCGA-KIRC at $n=131$ with $p^\star = 100, d^\star = 107$ in the companion
notebook.
\end{itemize}

\paragraph{Recommendations.} Choose $p^\star$ and $d^\star$ \emph{adaptively
from $n$}, not from the native modality dimensions:
\begin{itemize}
\item Target $p^\star + d^\star \le n/\tau$ with $\tau \ge 5$ as a working
rule (we use $p^\star, d^\star \le \lfloor n/\tau \rfloor$ per side in the
notebook; $\tau \in [5,10]$ trades bias for resolution). Reduce both
modalities by PCA when needed --- reducing only the larger side leaves the
estimator in regime~3.
\item Treat $\hat R_i^{\mathrm{cv}}$ (Section~\ref{sec:estimation-honest}) as
the primary estimand and report the in-sample $\hat R_i$ only with the gap
made explicit; the in-sample fraction of imaging variance attributable to
genomics should be reported alongside a cross-validated counterpart, e.g.\
$\sum_i \hat R_i^{\mathrm{cv}}/\min(p^\star, d^\star)$.
\item Power planning. The leading population correlation $\rho_1$ is detectable
above the Wachter bulk only when its squared value exceeds
$(\sqrt{\kappa_G(1-\kappa_X)} + \sqrt{\kappa_X(1-\kappa_G)})^2$; solving for $n$
gives a minimum sample size given a hypothesised effect. For
$\rho_1^2 = 0.2$ and a balanced $p^\star = d^\star$, this requires roughly
$n \gtrsim 25\,p^\star$, an order of magnitude beyond what cohort sizes in
radiogenomics usually offer --- reinforcing the need for aggressive
pre-reduction.
\item Probabilistic CCA (Section~\ref{sec:estimation-pcca}) with a small,
fixed latent dimension $k$ is the most parsimonious option when $n$ is the
binding constraint; rank-$k$ regularization sidesteps the high-dimensional
bulk by construction.
\end{itemize}

The practical message is conservative: in radiogenomic cohorts, sample size
is almost always the binding constraint and a defensible analysis is the one
whose effective $p^\star + d^\star$ is a small fraction of $n$, with both
sides reduced symmetrically and every reported number a cross-validated or
permutation-calibrated one.

Figure~\ref{fig:synthetic-sweep} illustrates these regimes on the synthetic
instance, where the planted truth is known. As $n$ grows with $p^\star, d^\star$
scaled to keep $\kappa_G + \kappa_X$ bounded, the in-sample $\hat\rho_1$ sits at
or near $1$ across the small-$n$ saturating regime --- indistinguishable from
the permutation null --- while the cross-validated $\hat R_1^{\mathrm{cv}}$ only
separates from zero once $n$ clears the Wachter edge. The empirical mutual
information shows the same threshold behaviour.

\begin{figure}[h]
\centering
\includegraphics[width=0.85\linewidth]{\synfig{sample_complexity_sweep.pdf}}
\caption{Sample-complexity sweep on the synthetic instance with known truth.
In-sample $\hat\rho_1$ tracks the null until $n$ exceeds the Wachter edge;
cross-validated $\hat R_1^{\mathrm{cv}}$ and empirical $\hat I(G;X)$ only
become informative past that threshold. Compare to
Figure~\ref{fig:kirc-sweep} for the same diagnostic on TCGA-KIRC.}
\label{fig:synthetic-sweep}
\end{figure}

\subsection{Honest recoverability: cross-validation and permutation}
\label{sec:estimation-honest}
In-sample canonical correlations are biased upward, severely so when $p, d$ are
not small relative to $n$. Two safeguards, both implemented in the notebook:
\begin{itemize}
\item \emph{Cross-validation.} Estimate $\hat a_i, \hat b_i$ on a training fold;
on the held-out fold form the scores $\hat a_i\T g$ and $\hat b_i\T x$ and report
their out-of-sample squared correlation as the honest $\hat R_i$. The gap to the
in-sample value quantifies overfitting.
\item \emph{Permutation test.} Repeatedly permute the patient labels of one
modality, refit, and collect the leading canonical correlation to build a null
distribution; the permutation $p$-value tests whether any genomic direction is
image-identifiable beyond chance.
\end{itemize}

Figure~\ref{fig:synthetic-cv} shows both diagnostics on the synthetic instance.
The in-sample $\hat R_i$ overshoots the planted truth by the amount predicted
by the high-dimensional bias, while the $5$-fold $\hat R_i^{\mathrm{cv}}$
tracks the truth for identifiable directions and collapses toward the null
$95\%$ quantile thereafter. Figure~\ref{fig:synthetic-perm} shows the
permutation null on $\hat\rho_1$: when the true signal is nonzero, the
observed value sits cleanly outside the null bulk and the per-direction
$p$-values fall below $0.01$ for the top three components, matching the
planted rank.

\begin{figure}[h]
\centering
\includegraphics[width=0.85\linewidth]{\synfig{cv_recoverability.pdf}}
\caption{Cross-validated recoverability on a synthetic instance with planted
rank $k=5$. In-sample $\hat R_i$ (grey) is biased upward; $5$-fold
$\hat R_i^{\mathrm{cv}}$ (blue) tracks the planted truth for the identifiable
directions and collapses to the null floor afterwards.}
\label{fig:synthetic-cv}
\end{figure}

\begin{figure}[h]
\centering
\includegraphics[width=0.75\linewidth]{\synfig{permutation_null.pdf}}
\caption{Permutation null on the leading canonical correlation
$\hat\rho_1$ for the synthetic instance. Red marks the observed value;
$p < 0.01$ for the top three directions, recovering the planted rank.}
\label{fig:synthetic-perm}
\end{figure}

The same optimism appears in the aggregate fraction of imaging variance
attributable to genomics. Figure~\ref{fig:synthetic-vardecomp} shows the
in-sample share $\sum_i \hat R_i / \min(p^\star, d^\star)$ alongside its
cross-validated counterpart on the synthetic instance: the in-sample number
sits well above the planted truth while the CV value tracks it. The same
diagnostic on TCGA-KIRC (Section~\ref{sec:kirc}) reports
$0.193$ in-sample against $0.012$ cross-validated --- a sixteen-fold gap that
is consistent with the synthetic mechanism rather than with any genuinely
recoverable signal.

\begin{figure}[h]
\centering
\includegraphics[width=0.75\linewidth]{\synfig{imaging_variance_decomposition.pdf}}
\caption{Share of imaging variance attributable to genomics on the synthetic
instance: in-sample (grey) versus $5$-fold cross-validated (blue), with the
planted truth marked. The in-sample share is biased upward by the
high-dimensional inflation described in
Section~\ref{sec:estimation-regimes}.}
\label{fig:synthetic-vardecomp}
\end{figure}

\section{Interpretation for radiogenomics}
\label{sec:interpretation}

The model reframes the radiogenomic question. Rather than testing whether CT
predicts each gene --- a high-dimensional, low-power multiple-testing problem ---
it asks for the \emph{recoverability spectrum} $\{\rho_i^2\}$ and the associated
genomic directions $\{a_i\}$. The outputs are: (1) a small number of
image-identifiable genomic axes, with gene/pathway loadings that make them
biologically legible; (2) a hard bound $\Sgcx$ on what any predictor can achieve,
per direction; (3) a total information budget $I(G;X)$; and (4) an honest,
cross-validated estimate of each. The governing principle is that radiogenomics
should target \emph{image-identifiable low-dimensional genomic subspaces}, whose
dimension is set by the biology--imaging channel and not by the gene panel size.

\section{A Bernoulli extension for mutation data via a Gaussian copula}
\label{sec:bernoulli}

The linear--Gaussian model of Section~\ref{sec:setup} is well matched to
continuous, approximately log-normal genomic readouts: log-CPM gene expression,
methylation $\beta$- or $M$-values, log-CNV ratios. It is plainly misspecified
for somatic mutation calls, where each $G_j \in \{0,1\}$ indicates the presence
of a non-silent variant in gene $j$ and the marginal frequencies $p_j$ are
small, gene-specific, and highly heterogeneous. This section extends the
framework to such data by replacing the Gaussian prior on $G$ with a
\emph{Gaussian-copula Bernoulli} (equivalently, multivariate probit) prior,
while preserving the canonical-correlation / mutual-information machinery of
Sections~\ref{sec:mi}--\ref{sec:canonical}.

\subsection{The latent-Gaussian generative model}
\label{sec:bern-model}

Let $p_j = \Pr(G_j = 1)$ and $\tau_j = \Phi^{-1}(1-p_j)$, where $\Phi$ is the
standard normal CDF. Introduce a latent vector
\[
  Z \sim \mathcal N(0, \Sigma_Z), \qquad
  \mathrm{diag}(\Sigma_Z) = I_p,
\]
and define the observed mutation indicator by thresholding:
\[
  G_j = \mathbf{1}\!\{Z_j > \tau_j\}, \qquad j = 1,\ldots,p.
\]
By construction the marginals are correct, $G_j \sim \mathrm{Bern}(p_j)$, and
the pairwise joint laws are governed by the tetrachoric correlation $\Sigma_{Z,jk}$,
which can be any positive-definite correlation matrix --- providing arbitrary
dependence structure between mutated genes (e.g.\ mutual exclusivity,
co-occurrence within pathways).

Imaging is coupled to the \emph{latent} state, not to the indicators directly:
\begin{equation}
  X \mid Z \sim \mathcal N\!\bigl(B\,Z,\ \sigma^2 I_d\bigr),
  \qquad
  B \in \R^{d \times p}.
  \label{eq:bern-channel}
\end{equation}
This is the natural specification: the biological mechanism that produces an
imageable phenotype (proliferation, vascularity, necrosis) responds to an
underlying continuous burden, of which an observed mutation is a coarse,
binarized readout. The joint structure factors as
\[
  Z \to G,\qquad Z \to X,
\]
so $G$ and $X$ are conditionally independent given $Z$.

\subsection{Mutual information and a tight DPI bound}
\label{sec:bern-mi}

Because $G$ is a deterministic, memoryless function of $Z$, the data
processing inequality yields the upper bound
\begin{equation}
  I(G;X) \;\le\; I(Z;X),
  \label{eq:bern-dpi}
\end{equation}
with equality only if $Z$ can be recovered from $G$ (impossible whenever any
$p_j \in (0,1)$, since thresholding is lossy). The right-hand side reduces
exactly to the Gaussian formula of Section~\ref{sec:mi}: if $\{\rho_i\}_{i=1}^{k}$
denote the canonical correlations of $(Z,X)$ under \eqref{eq:bern-channel}, then
\begin{equation}
  I(Z;X) \;=\; -\tfrac{1}{2}\sum_{i=1}^{k} \log\bigl(1-\rho_i^2\bigr),
  \qquad
  k = \rank\bigl(\Sigma_Z^{1/2} B^\top \Sigma_X^{-1} B \Sigma_Z^{1/2}\bigr),
  \label{eq:bern-mi}
\end{equation}
with $\Sigma_X = B \Sigma_Z B^\top + \sigma^2 I_d$. The recoverability spectrum,
the directions $\{a_i\}$, and the eigenproblem of
Section~\ref{sec:eigenproblem} all transfer verbatim, with $\Sigma_G$ replaced
by $\Sigma_Z$.

\paragraph{The binarization gap.}
The slack in \eqref{eq:bern-dpi} is
\[
  I(Z;X) - I(G;X) \;=\; H(Z\mid G) - H(Z\mid G,X) \;=\; I(Z;X \mid G),
\]
the residual information about $Z$ not already revealed by the mutation call.
For a single gene with frequency $p_j$, the entropy lost to binarization is
$H(Z_j) - H(G_j) = \tfrac12 \log(2\pi e) - h_2(p_j)$ where $h_2$ is the binary
entropy in nats; rare mutations ($p_j \ll 1/2$) are highly lossy, common ones
near $p_j = 1/2$ less so. Thus \eqref{eq:bern-dpi} is tightest for cohorts
dominated by intermediate-frequency mutations and loosest for rare-variant
panels --- a useful guide to when the bound is informative.

\subsection{Per-gene recoverability and AUC ceilings}
\label{sec:bern-pergene}

For a named mutation $G_j$, the Bayes-optimal predictor from $X$ is the
posterior probability
\[
  \pi_j(x) \;=\; \Pr(G_j = 1 \mid X = x) \;=\; \Phi\!\left(\frac{\mu_j(x) - \tau_j}{s_j}\right),
\]
where $(\mu_j(x), s_j^2) = (\E[Z_j \mid X=x],\ \Var(Z_j \mid X=x))$ are the
linear-Gaussian posterior moments for the latent coordinate, computed from
$(\Sigma_Z, B, \sigma^2)$ exactly as in Section~\ref{sec:posterior}. The
\emph{latent recoverability} $R_j^Z = \Var(\E[Z_j\mid X]) / \Var(Z_j)$ is the
direct analogue of $R(v)$ in the Gaussian case, and the binary AUC of any
predictor $\hat G_j(x)$ is bounded by
\begin{equation}
  \mathrm{AUC}_j \;\le\; \Phi\!\left(\sqrt{\tfrac{R_j^Z}{1-R_j^Z}}\Big/\sqrt{2}\right),
  \label{eq:bern-auc}
\end{equation}
attained by the Bayes classifier $\pi_j(x) > p_j$. Equation~\eqref{eq:bern-auc}
is the natural object to compare against the empirical per-gene
$\mathrm{AUC}_{\mathrm{LR}}$ and $\mathrm{AUC}_{\mathrm{RF}}$ of
Section~\ref{sec:kirc-pergene}: an empirical AUC near $0.5$ is consistent with
either a small $R_j^Z$ \emph{or} a degenerate $p_j$, and \eqref{eq:bern-auc}
disentangles the two.

\subsection{Estimation}
\label{sec:bern-estimation}

The estimator splits into three independent pieces, each well-studied:

\begin{enumerate}
\item[(i)] \textbf{Marginal thresholds.} Estimate $\hat p_j = \bar G_{\cdot j}$
and set $\hat\tau_j = \Phi^{-1}(1-\hat p_j)$. Apply a continuity correction or
empirical-Bayes shrinkage for rare mutations.
\item[(ii)] \textbf{Latent correlation $\hat\Sigma_Z$.} Estimate pairwise
tetrachoric correlations from $2{\times}2$ tables of $(G_j, G_k)$ by maximum
likelihood under the bivariate-normal copula, then project the resulting matrix
to the nearest positive-definite correlation matrix (e.g.\ Higham's
\texttt{nearcorr}). For high-dimensional panels apply Ledoit--Wolf shrinkage as
in Section~\ref{sec:estimation-cov}.
\item[(iii)] \textbf{Cross-modal covariance $\hat\Sigma_{ZX}$.} The
point-biserial covariance $\Cov(G_j, X_\ell)$ is related to the latent
covariance by
\[
  \Cov(G_j, X_\ell) \;=\; \phi(\tau_j)\,\Cov(Z_j, X_\ell),
\]
where $\phi$ is the standard normal density. Estimate $\Cov(G_j,X_\ell)$ from
the sample and rescale by $\phi(\hat\tau_j)$ to obtain $\hat\Sigma_{ZX, j\ell}$.
\end{enumerate}

With $(\hat\Sigma_Z, \hat\Sigma_{ZX}, \hat\Sigma_X)$ in hand, the regularized
CCA of Section~\ref{sec:estimation-cca}, the permutation null of
Section~\ref{sec:estimation-honest}, and the empirical MI bound carry over
without further change. The dimensional-collapse warnings of
Section~\ref{sec:estimation-regimes} are if anything \emph{more} acute, since
tetrachoric estimation has higher variance than Pearson correlation,
particularly for rare-rare gene pairs.

\subsection{Mixed-modality genomics}
\label{sec:bern-mixed}

In practice a single cohort carries both continuous (expression, CNV) and
binary (mutation) features. The copula formulation accommodates this by
partitioning the latent vector as $Z = (Z^{\mathrm{cts}}, Z^{\mathrm{bin}})$
with a single joint $\Sigma_Z$. The continuous block uses the identity link
($G^{\mathrm{cts}} = Z^{\mathrm{cts}}$ up to marginal standardization) and the
binary block uses the probit link of Section~\ref{sec:bern-model}. Estimation
of $\Sigma_Z$ proceeds blockwise --- Pearson within the continuous block,
tetrachoric within the binary block, polyserial across blocks --- after which
the rest of the pipeline is identical. This is the route by which the present
manuscript's gene-expression analysis (Section~\ref{sec:tcga-kirc}) and the
per-gene mutation validation (Section~\ref{sec:kirc-pergene}) can be unified
into a single coherent recoverability calculation.

\section{Application: TCGA-KIRC}
\label{sec:kirc}
\label{sec:tcga-kirc}

We now instantiate the framework on a real radiogenomic cohort --- TCGA-KIRC
(clear-cell renal cell carcinoma) --- to demonstrate how the recoverability
spectrum, the cross-validated and permutation-calibrated diagnostics of
Section~\ref{sec:estimation-honest}, and the dimensional-collapse warnings of
Section~\ref{sec:estimation-regimes} interact on a typical cohort. The point
of this section is \emph{not} to advertise a predictor: the result is a null
one. The point is to convert that null into a quantitative empirical ceiling
on the imaging-to-genomics channel and a sample-size design recommendation,
exactly as the synthetic experiments of
Sections~\ref{sec:canonical}--\ref{sec:estimation-honest} predict. All numbers
below are reproducibly generated by the companion notebook and saved to
\texttt{notebooks/figures/tcga\_kirc/gene\_expression/tumor\_radimagenet/%
stats.json}.

\subsection{Data and processing pipeline}
\label{sec:kirc-pipeline}
The two modality \texttt{AnnData} objects fed to the notebook are produced by
two scripts in \texttt{scripts/}.

\paragraph{Imaging
(\texttt{process\_imaging\_tcga\_kirc.py} $\to$ \texttt{make\_imaging\_matrix.py}).}
For each TCGA-KIRC patient we obtain the venous-phase abdominal CT volume
(\texttt{0502\_VENOUS.nii}) and a manually curated tumour mask. The pipeline:
\begin{enumerate}
\item Pull the canonical TCIA volumes and the radiologist-drawn tumour
segmentations; orient every volume into the canonical
nibabel orientation.
\item Run TotalSegmentator with the
\texttt{kidney\_left}/\texttt{kidney\_right} target list to produce an organ
mask; combine with the tumour mask to form (tumour, organ, tumour$\cup$organ)
mask triples per patient. Small blobs are removed, holes filled, and a
morphological closing applied to stabilise mask boundaries.
\item Extract patient-level imaging features. We use a pretrained
\emph{RadImageNet} 2D backbone applied tile-wise within the tumour mask
(\texttt{rgit.radimagenet.get\_radimagenet\_embeddings}), producing a
$d_{\mathrm{native}} = 1841$-dimensional embedding per patient. (A second
pipeline using PyRadiomics shape/texture features is supported for
ablations.)
\item Stack patient-level embeddings into an \texttt{AnnData} object whose
\texttt{obs} index is the TCGA \texttt{patient\_id} and whose \texttt{var}
index is the embedding-dimension labels. Output:
\texttt{data/tcga\_kirc/imaging/tumor\_radimagenet.h5ad}.
\end{enumerate}

\paragraph{Genomics (\texttt{make\_genomics\_matrix.py}).}
For each matched patient we pull the TCGA-KIRC RNA-seq counts (raw STAR counts
from the GDC) and the somatic-mutation MAFs.
\begin{enumerate}
\item Counts are CPM-normalised, $\log_{1+}$-transformed, and reduced to the
top $2000$ highly variable genes (Seurat-style HVG selection in
\texttt{scanpy.pp.highly\_variable\_genes}); the resulting matrix is stored
in the primary \texttt{AnnData}.
\item Somatic mutations are folded into a sparse binary
gene-symbol matrix on the same patient index for downstream per-gene
validation (Section~\ref{sec:kirc-pergene}).
\item Output: \texttt{data/tcga\_kirc/genomics/gene\_expression.h5ad}.
\end{enumerate}
After intersecting patients across both \texttt{AnnData} objects we retain
$n = 190$ patients with $p_{\mathrm{native}} = 2000$ HVG transcripts and
$d_{\mathrm{native}} = 1841$ RadImageNet features.

\subsection{Working-dimension budget}
\label{sec:kirc-cohort}
Both modalities are $z$-scored per feature and PCA-reduced symmetrically to
working dimensions
$p^\star = d^\star = 38 = \lfloor n / 5 \rfloor$ following the
$p^\star + d^\star \le n/\tau$ rule of
Section~\ref{sec:estimation-regimes} ($\tau = 5$). This places the analysis at
$\kappa_G + \kappa_X = 76/190 \approx 0.4$ --- the high-dimensional regime,
not the saturating one, but well inside the territory where in-sample
canonical correlations are heavily biased upward.

\subsection{The recoverability spectrum and its null}
\label{sec:kirc-spectrum}
Fitting regularised CCA (Ledoit--Wolf shrinkage on both sides) returns an
in-sample spectrum dominated by a near-unit leading correlation
($\hat\rho_1 \approx 1.00$, $\hat R_1 \approx 1.00$) and a total in-sample
mutual information of $21.3$ bits over $12$ retained components
(Fig.~\ref{fig:kirc-spectrum}). This is exactly the geometric artefact
described in regime~2--3 of Section~\ref{sec:estimation-regimes}:
$5$-fold cross-validation collapses $\hat R_1$ from $\approx 1$ to
$\hat R_1^{\mathrm{cv}} = 0.094$, and held-out prediction of the leading
canonical genomic score has $R^2 = -3.74$ --- worse than predicting the mean.
The companion permutation test ($N_{\mathrm{perm}} = 200$) returns
$p$-values $0.32, 0.16, 0.24$ for the top three canonical correlations.
None of the leading directions are distinguishable from chance.

\begin{figure}[h]
\centering
\includegraphics[width=0.85\linewidth]{\kircfig{cv_recoverability.pdf}}
\caption{Recoverability spectrum on TCGA-KIRC. In-sample $\hat R_i$ (grey)
tracks the high-dimensional bulk; cross-validated $\hat R_i$ (blue, $5$-fold)
collapses to a flat residual near zero. The gap is the overfitting predicted
in Section~\ref{sec:estimation-regimes}, regime~2--3. Contrast with the
synthetic instance of Fig.~\ref{fig:synthetic-cv}, where CV tracks the
planted truth.}
\label{fig:kirc-spectrum}
\end{figure}

\subsection{An empirical upper bound on $I(G;X)$}
\label{sec:kirc-ucb}
Section~\ref{sec:estimation-regimes} predicts that, in this regime, the
permutation null on $\hat\rho_1$ should itself be pulled toward $1$ by the
high-dimensional geometry; this is borne out: the per-permutation
$\hat\rho_1^{\mathrm{null}}$ has mean $0.84$ and a $95\%$ quantile of $1.00$.
We therefore convert the test into a one-sided upper confidence limit by
aggregating, across permutations, the cumulative mutual information
$\mathrm{MI}^{\mathrm{null}}(K) = -\tfrac12 \sum_{i\le K}
\log(1 - \hat\rho_{i,\mathrm{null}}^2)$ over the top $K=3$ tested components,
and reading off its $95$th percentile (Fig.~\ref{fig:kirc-ucb}):
\begin{equation}
I(G;X)\ \big|_{\text{KIRC, top-}3}\ \le\ 21.2\ \mathrm{bits}\quad
(\text{one-sided }95\%\text{ UCL}),
\label{eq:kirc-ucl}
\end{equation}
with observed cumulative MI $20.5$ bits sitting close to the null mean of
$8.0$ bits expressed in bits ($5.5$ nats). The bound is wide --- this is the
honest finding, not a weakness of the method --- because the null is itself
saturating; tightening it requires either larger $n$ or smaller
$p^\star, d^\star$.

\begin{figure}[h]
\centering
\includegraphics[width=0.75\linewidth]{\kircfig{mi_upper_bound.pdf}}
\caption{Cumulative-MI null distribution (top-$3$ canonical components) on
TCGA-KIRC. The observed value (red) lies within the null bulk; the dashed
line marks the $95\%$ upper confidence limit on $I(G;X)$ used
in~\eqref{eq:kirc-ucl}.}
\label{fig:kirc-ucb}
\end{figure}

A more useful summary is the \emph{effective image-identifiable rank}: the
number of canonical components whose cross-validated $\hat R_i^{\mathrm{cv}}$
exceeds the per-component permutation-null $95\%$ quantile. On this cohort,
\begin{equation}
\hat k_{\mathrm{eff}} = 0\ \text{(of }K = 3\text{ tested),}
\end{equation}
so the empirical instantiation of the rank bound of
Proposition~\ref{prop:rank}, after sample-size deflation, is zero.

\subsection{Sample-complexity sweep}
\label{sec:kirc-sweep}
A single-$n$ null result is most informative when paired with the scaling of
that null in $n$. We subsample the cohort to
$n \in \{50, 57, 95, 133, 162, 190\}$, refit the working pipeline (same
$p^\star + d^\star \le n/\tau$ budget, $5$-fold CV, $100$ permutations) and
plot the three observables vs.\ $n$ (Fig.~\ref{fig:kirc-sweep}). The
cross-validated $\hat R_1^{\mathrm{cv}}$ is flat at $\approx 0.05$--$0.10$
across the full range; doubling $n$ does not move it. The Wachter detection
edge sits at $\approx 0.63$ throughout, because $p^\star, d^\star$ grow with
$n$ under our rule, so the cohort-statistical ceiling for distinguishing
$\rho_1$ from the high-dimensional bulk would require
$\rho_1^2 \gtrsim 0.6$ --- well above anything observed at any subsample
size. Inverting the Wachter expression for a target $\rho_1^2 = 0.2$ at
balanced $p^\star = d^\star$ gives the sample-size requirement
$n \gtrsim 25\,p^\star$ quoted in Section~\ref{sec:estimation-regimes}.

\begin{figure}[h]
\centering
\includegraphics[width=0.85\linewidth]{\kircfig{sample_complexity_sweep.pdf}}
\caption{Sample-complexity sweep on TCGA-KIRC. Honest $\hat R_1^{\mathrm{cv}}$
(blue squares) is flat across all subsampled cohort sizes;
in-sample $\hat\rho_1^2$ (grey) is dominated by the high-dimensional bulk;
the permutation $95\%$ null (red triangles) and Wachter edge (green crosses)
sit above the observed honest recoverability at every $n$.}
\label{fig:kirc-sweep}
\end{figure}

\subsection{Per-gene non-linear validation}
\label{sec:kirc-pergene}
The recoverability spectrum makes a sharp prediction: genes loading on
non-identifiable directions should not be predictable from imaging, and
flexible non-linear models should not substantially exceed the linear
Bayes-optimal ceiling. Picking five high- and five low-recoverability
mutated genes (mutation frequency in $[0.05, 0.95]$) and training $L_2$
logistic regression and a small random forest to predict mutation status
from the imaging PCs ($5$-fold CV), we observe
\begin{center}
\begin{tabular}{lcc}
\toprule
group               & mean AUC (logistic) & mean AUC (RF) \\
\midrule
image-identifiable  & $0.499$ & $0.522$ \\
non-identifiable    & $0.497$ & $0.490$ \\
\bottomrule
\end{tabular}
\end{center}
Both groups sit at chance and the non-linear model does \emph{not} exceed
the linear ceiling --- consistent with the linear--Gaussian model and with
the global null finding (Fig.~\ref{fig:kirc-pergene}).

\begin{figure}[h]
\centering
\includegraphics[width=0.85\linewidth]{\kircfig{per_gene_auc.pdf}}
\caption{Per-gene mutation predictability on TCGA-KIRC. Five
high-recoverability and five low-recoverability mutated genes, two
classifiers, $5$-fold CV. All bars sit within sampling noise of the
chance line at $0.5$.}
\label{fig:kirc-pergene}
\end{figure}

\subsection{Implications}
\label{sec:kirc-implications}
Three concrete take-aways follow.
\begin{enumerate}
\item \emph{Empirical ceiling.} For TCGA-KIRC at $n = 190$ with
$2000$ HVG-filtered transcripts and $1841$ RadImageNet features, the
imaging-to-genomics channel admits a $95\%$ upper bound
$I(G;X) \le 21.2$ bits over the leading three canonical directions
(Eq.~\ref{eq:kirc-ucl}), with effective image-identifiable rank
$\hat k_{\mathrm{eff}} = 0$. The substantive empirical claim is that no
genomic direction in this cohort is image-identifiable above chance under
the working-dimension budget $\tau = 5$.
\item \emph{Design recommendation.} The sample-complexity sweep
(Section~\ref{sec:kirc-sweep}) shows $\hat R_1^{\mathrm{cv}}$ stationary in
$n$ over the range $[50, 190]$; the Wachter edge analysis indicates that
detecting $\rho_1^2 = 0.2$ at a balanced working dimension
$p^\star = d^\star$ requires $n \gtrsim 25\,p^\star$, which under our
adaptive rule is an order of magnitude beyond typical radiogenomic cohort
sizes. The framework converts the null result into a quantitative cohort
specification.
\item \emph{Audit of literature predictors.} Published radiogenomic AUCs
in the range $0.65$--$0.80$ at cohort sizes $n \sim 100$--$300$ imply
per-gene recoverabilities meaningfully above the cohort-statistical
ceiling derived above for KIRC. The non-linear validation of
Section~\ref{sec:kirc-pergene} matches but does not exceed the linear
Bayes-optimal ceiling --- consistent with the model. Any study reporting
substantially higher per-gene AUCs at comparable cohort dimension should
be treated as a candidate for overfit or leakage rather than as evidence
against the framework.
\end{enumerate}

\appendix
\section{Symbol reference}
\label{app:symbols}

\noindent
\begin{tabular}{ll}
$g \in \R^p,\ x \in \R^d$ & per-patient genomic / imaging feature vectors\\
$\mathbf{G} \in \R^{n\times p},\ \mathbf{X} \in \R^{n\times d}$ & stacked data matrices ($n$ patients)\\
$A \in \R^{d\times p}$ & genomics-to-imaging transfer map\\
$\sigma^2$ & isotropic imaging noise variance\\
$\Sg, \Sx, \Sgx$ & genomic / imaging / cross covariance\\
$\Sgcx$ & posterior covariance of $G$ given $X$ (recoverability floor)\\
$B = \Sgx \Sx^{-1}$ & Bayes-optimal linear map $x \mapsto \E[g\mid x]$\\
$M = \sigma^{-1} A \Sg^{1/2}$ & whitened transfer operator; $d_i$ its singular values\\
$\rho_i,\ R_i = \rho_i^2$ & $i$-th canonical correlation / recoverability\\
$a_i,\ b_i$ & $i$-th canonical genomic / imaging direction\\
$R(v)$ & recoverability score of a named genomic direction $v$\\
$I(G;X)$ & mutual information between modalities\\
$k$ & dimension of the image-identifiable genomic subspace\\
\end{tabular}

\end{document}
